\documentclass[12pt]{article}

\begin{document}

\textbf{CHAPTER 2
BASIC CONCEPTS IN 2D}

\hfill 

In the preceding chapter we introduced a brief list of basic concepts
of discrete dynamics. Here, we expand on these concepts in
the one-dimensional context, in which, uniquely, we have the
advantage of a simple graphical representation. The official,
abstract definitions of all these concepts may be found in the
Appendices.

\hfill

\textbf{1.1 MAPS}

\hfill

By a map we mean a continuous function from a space, called
the domain, to itself. 1 In the one-dimensional context, the domain
might be an interval (with or without endpoints) of the real number
line, or even the entire line.

\hfill

If f is a map on a real interval, I, we indicate this in symbols
$f:I \to I$. For example, if the map is defined by the rule $f(x) = x^2$,
and I is the closed interval [-2, 2], we may visualize the map graphically,
as shown in Figure 2-1. The action of the map is to move
points from the horizontal axis to the vertical axis in two strokes:

\begin{itemize}
\item vertically from the horizontal axis to the graph,
\item horizontally from the graph to the vertical axis,
\end{itemize}

as shown in Figure 2-2.

\hfill


The image (or range) of the map is the set of all points obtained
asf(x) while x takes on all values in the domain, I, and is written in
symbols as f[I]. For a point y in I, a pre image of rank 1 is a point x 
in I that is mapped to y; that is, x is a preimage of rank 1 of y if y =
f(x). A preimage of rank 2 of y is a preimage of rank 1 of a preimage
of rank 1, and so on. Every point y in I has a set of preimages of
every rank, which may be empty. Determining all preimages of a
point creates a genealogical tree, called the arborescent sequence of
preimages.

\hfill

The map f is one-to-one, if, for every point y in I, the set of preimages
of y has either no points or just one point. For example, the
map f defined by the same rule, $f(x) = x^2$ , but with the smaller
domain [0,1], is one-to-one. A one-to-one map has a unique inverse
map, $f^{-1}: F[I] \to I$ which undoes what f does. This can be visualized
on the graph off as a motion in two strokes:

\begin{itemize}
\item horizontally from the vertical axis to the graph off,
\item vertically from the graph to the horizontal axis,
\end{itemize}

as shown in Figure 2-3.

Note that if we try to invert the map of Figure 2-2 by this twostroke
process, we discover all of the preimages of a given point y
in I (represented as the vertical axis) in one step. This is shown in
Figure 2-4, where we find two preimages.

\hfill

In this text we will be concerned exclusively with the lowest
step of the staircase to chaos, the ID case. We will be interested
especially with maps which are not one-to-one. These are called
many-to-one, or noninvertible, maps. For such maps, points generally
have more than one preimage of rank 1, and the number of
preimages of a given rank determines a zone of multiplicity, discussed
below.

\hfill

\textbf{2.2 MULTIPLICITIES}

\hfill

Given any map of an interval I, we may choose a point y in I on
the vertical axis, locate all preimages by the graphical method, and
count them up. Thus, we may decompose the vertical axis into sets
of points all sharing the same number of preimages. We denote
these zones by:

\begin{itemize}
\item Zo (all points having no preimages)
\item ZI (all points having exactly one preimage)
\item Z2 (all points having two distinct preimages)
\end{itemize}

and so on.

\hfill

These are called the layer sets of the map; those sets that are
nonempty and open (that is, have no boundary points) are called
multiplicity zones. The zones $Z_o$ and $Z_2$ are shown in Figure 2-5.
They exhaust the whole interval I = [-2, 2] except for the point 0,
which separates the two zones. We also say this map is of type $Z_o$ -
$Z_2$, meaning there are two zones, one of multiplicity zero, the other
of multiplicity two. The suffix indicates the multiplicity. The interval
Z2 , may be considered a folded image, that is, two halves of the
domain I are folded onto this image. For each half of the domain,
our map does have an inverse. These are called partial inverses.

\hfill

The point 0 is the only point of $Z_1$ in this example; that is, it has
multiplicity one (its unique preimage is 0). It is called a critical
point because it lies on the boundary of two zones. In fact, we can
describe this map as a nonlinear folding. That is, the map folds the
horizontal axis at the critical point, then stretches them in a nonlinear
fashion onto the range interval. This is why the critical point is
sometimes called afold point.

\hfill

Generally, we will be interested in relatively simple maps, such
as polynomials, in which only finite multiplicities, with generic
(that is, typical) fold points, are encountered. We call these finitely
folded maps. For example, a typical cubic map has multiplicities 1
and 3, and we say it is of type $Z_1 - Z_3 - Z_1$. These basic concepts of
iteration theory should be approached through a careful study of
simple (for example, polynomial) examples.

\hfill

\textbf{2.3 TRAJECTORIES AND ORBITS}

\hfill

A map generates a discrete dynamical system by iteration. That
is, the map is applied again and again, and points move along a dotted
path called a trajectory. For example, choosing an initial point
$x_0$, let $x_1$ denote its image under the map f, likewise $X_2$ the image of $x_1$ and so on. The infinite sequence $(x_0, x_1, x_2, ...)$ is the trajectory
of $x_0$. This sequence may jump around a finite set of points. The
minimum set of points which holds a trajectory its called its orbit.
When finite, an orbit is called cyclic, or periodic, and the number of
its points is its order. A fixed point, defined by f(x) = x, is a special
kind of periodic point, the orbit of which is a single point. It has
order 1. If an orbit contains only two points, it is called a 2-cycle,
and so on.

\hfill

We now describe a graphical method for plotting trajectories,
called the Koenigs-Lemeray method. Note that in our graphs, both
axes represent the same set, since the domain and range of our map
consist of the same interval.

\hfill

Given $x_0$, envisioned on the horizontal axis, we may find Xl on
the vertical axis by the two-stroke method described in 2.1. Next,
we must repeat this process, starting from the point $X_1$ on the horizontal
axis. Our immediate problem, then, is to transfer the distance
Xl from the vertical axis to the corresponding distance on the horizontal
axis.

\hfill

One way to carry out this transfer is shown in Figure 2-6. Here
we use a compass to measure the vertical distance, Xl' and rotate it
to the horizontal distance, $x_1$. Another method is to use a protractor
to construct a line descending at slope -1, or 45 degrees, as shown
in Figure 2-7.

\hfill

Yet another method - and this is the one we prefer - is shown
in Figure 2-8. We draw a line from the lower left comer, ascending
at slope 1. This line is called the diagonal (in symbols, $\Delta$). Now,
using only a square, we draw a horizontal line from $x_1$ on the vertical
axis until it meets Il., then draw a vertical line until it meets the
horizontal axis. This determines the horizontal distance, $x_1$ as
shown in Figure 2-8.

\hfill

The entire construction from horizontal $x_0$ to vertical $x_1$ to horizontal $x_1$ may now be summarized as follows:

\begin{itemize}
\item vertical from horizontal axis to graph,
\item horizontal from graph to vertical axis,
\item horizontal from vertical axis to diagonal,
\item vertical from diagonal to horizontal axis.
\end{itemize}

This construction may be abbreviated somewhat since the third
stroke retraces (undoes) part of the second, as shown in the four
strokes of Figure 2-9. The abbreviated construction (Figure 2-10) is:

\begin{itemize}
\item vertical from horizontal axis to graph,
\item horizontal from graph to diagonal,
\item vertical from diagonal to horizontal axis.
\end{itemize}

This is the three-stroke graphical method for plotting one step
of a trajectory within the horizontal axis. When we proceed to plot
the point $x_2$ by this method, however, we find a further opportunity
for abbreviation. The last stroke above, which locates $x_1$ on the horizontal
axis, i.e.,

\begin{itemize}
\item vertical from diagonal to horizontal axis,
\end{itemize}

is followed by the first stroke of the second step,

\begin{itemize}
\item vertical from horizontal axis to graph,
\end{itemize}

which may be combined into a single stroke,

\begin{itemize}
\item vertical from diagonal to graph.
\end{itemize}

Thus the iterated sequence, beginning with Xo on the horizontal
axis, is:

\begin{itemize}
\item vertical from horizontal axis to graph,
\item horizontal from graph to diagonal,
\item vertical from diagonal to graph,
\end{itemize}

and continue. After the first step, we may pretend that we are jumping
about on the diagonal, which after all is just another copy of the
domain interval I. Each step has two strokes,

\begin{itemize}
\item vertical to the graph; horizontal to the diagonal
\end{itemize}

(which we may remember by the mnemonic, vertigo-horrid), as
shown in Figure 2-11. That is the graphical method of KoenigsLemeray,
also known as the staircase method, or cobweb construction.
Using it, we may quickly follow trajectories for several jumps
on the diagonal. See Figure 2-12.

\hfill

Using this method, one may graphically verify this useful fact:
if a point returns to its starting point after two iterations of the map,
the starting point is either a fixed point or a 2-periodic point of the
map.

\hfill

\textbf{2.4 ATTRACTORS, BASINS,AND BOUNDARIES}

\hfill

Try out the staircase method using the Myrberg map, 1
$f(x) = x^2 -c$ ,with the entire real line as the domain, and various
values for the control parameter, c. You may quickly find that some
trajectories converge to afixed point, while others run off to positive
infinity (upper right) on the diagonal. The fixed points are seen
immediately as the crossing points of the graph and the diagonal,
and are defined by the property: f(x) = x.

\hfill

An interval is called trapping if it is mapped into itself, and
invariant if it is mapped exactly onto itself. If a bounded interval is
trapping, then all of its trajectories are trapped inside, and must
converge to a closed, invariant, and bounded limit set. These limit
sets are the attractors of the map. Attractors may be classified in
three categories:

\begin{itemize}
\item a point attractor is a single point,
\item a cyclic attractor is a finite set of points, and
\item a chaotic attractor is any other type of attractor.2
\end{itemize}

The basin of an attractor is the set of all points tending to that
attractor. The domain is decomposed into the basins of different
attractors, including the basin of infinity, which consists of all
points whose trajectories run away from any bounded set.

\hfill

The boundaries of the basins, also called frontiers or separatrices,
are of primary importance in dynamical systems theory. A
detailed study of a map results in a portrait, in which the domain is
decomposed into basins, one attractor shown in each.

\begin{enumerate}
\item Myrberg was one of the first to study the bifurcation sequence of this map. See the Bibliography
for references to his work.
\item This reflects the fact that different definitions of chaos abound in the literature.
\end{enumerate}

\hfill

\textbf{2.5 BIFURCATIONS}

\hfill

As in the Myrberg map, $f(x) = x^2 - c$, we frequently encounter
maps which depend on a parameter. As the parameter is changed,
the portrait of the attractive set of the map may change gradually
and insignificantly; however, as certain special values of the parameter
are crossed, there may be a sudden and significant change in
the portrait of the map. These special values are called bifurcation
points, and the sudden changes in the portrait are called bifurcations.
At the present time, dynamical systems theory does not have
a satisfactory and rigorous definition of bifurcation, but the subject
is now evolving through the study of examples. In fact, the goal of
this book is to describe some of these examples, in a two-dimensional
context.

\hfill

In the current one-dimensional context, we again have the benefit
of an excellent visualization device, the response diagram. In the
case of a single control parameter, this is a two-dimensional graphic
in which the vertical axis represents the domain of the map and the
horizontal axis represents the control parameter. Above each point
on the horizontal axis, the portrait of the corresponding map is indicated,
with its attractors, basins, and basin boundaries. For a oneparameter
family of maps of a two-dimensional domain, the
response diagram is three-dimensional, as we will soon see.

\hfill

\textbf{2.6 EXEMPLARY BIFURCATION}

\hfill

The simplest bifurcations are the fold and the flip. These may
involve changes to any kind of attractor. To introduce the basic concepts
of bifurcation theory, however, we will describe the fold
bifurcation in the simplest case, which involves point attractors.

\hfill

The fold bifurcation is a catastrophic bifurcation. This means
that, as the control parameter varies, an attractor appears or disappears
suddenly. In this event, as shown in Figure 2-13 with the
control parameter moving to the right on the horizontal axis, a fixed
point appears, and immediately separates into a pair of distinct
fixed points. One is an attractor, the other, a repellor. The repellor is
shown below the attractor. Points between the two fixed points are
attracted to the upper fixed point, and repelled by the lower fixed
point. These tendencies are indicated by the arrows in Figures 2-15
to 2-17.

\hfill

To understand the mechanism of this bifurcation, we now turn
to a specific example, the Myrberg family of maps, f(x) = x2 -c.
The graph of a map of this family is an upward-opening parabola,
with the vertex on the vertical axis at distance c below the horizontal
axis. As c increases, the parabola moves downward. Three cases
of this graph are shown in Figs. 2-14, 2-15, and 2-16.

\hfill

In the first case, Figure 2-14, with c = - 0.5, the parabola does
not meet the diagonal because for this value of c, there are no fixed
points. All trajectories tend upward without bound, to infinity.

\hfill

In the next case, Figure 2-15, with c = - 0.25, the parabola
meets the diagonal in a single point, which is the fixed point x = 0.5,
corresponding to this value of c, the bifurcation value. Trajectories
approach from below, but depart from above.

\hfill

In the last case, Figure 2-16, with c = 0, the parabola cuts the
diagonal in two points, the fixed points x = 0 and x = I, which are,
respectively, an attractor and a repellor.

\hfill

The flip is a subtle bifurcation. This means that, in contrast to
catastrophic and explosive bifurcations, its effect is too subtle to
observe at the moment of bifurcation when the control parameter
passes its critical value, but becomes apparent later, as the parameter
continues to increase. In the flip, a point attractor loses its
attractiveness. From it is emitted a cyclic attractor of period 2.

\hfill

A response diagram of this event is shown in Figure 2-17. To
the left of the bifurcation value of the control parameter (horizontal
axis) there is a single fixed point, and it is an attractor, FP+. The
attraction in the vertical state space is shown by the heavy arrows.
To the right, there is still only one fixed point, but it is a repellor,
FP-. But there is also a 2-cycle, which is attractive, 2P+. Looking
only at the attractors in the picture, we see that the attractive point
has been replaced by an attractive 2-cycle, as the control parameter
moves to the right. At first, the two points of this 2-cycle are very
close together, then they gradually separate.

We now move on to two dimensions.

\begin{thebibliography}{9}
\bibitem{texbook}
RALPH H. ABRAHAM, LAURA GARDIINI, CHRISTIAN MIIRA, (2016) \emph{Chaos Discrete Dynamical Systems}, Springer.

\end{thebibliography}

\end{document}

